\documentclass[a4paper,11pt]{article}

\usepackage[utf8]{inputenc}
\usepackage[french]{babel}
\usepackage{amsmath}
\usepackage{graphicx}
\usepackage{float}
\usepackage{xcolor,listings}
\usepackage{geometry}

\geometry{
 a4paper,
 total={170mm,257mm},
 left=20mm,
 top=20mm,
 }

\lstset{
  basicstyle=\ttfamily\,
  columns=fullflexible,
  frame=single,
  breaklines=true,
  postbreak=\mbox{\textcolor{red}{$\hookrightarrow$}\space},
  showstringspaces=false,
}

\definecolor{main-color}{rgb}{0.6627, 0.7176, 0.7764}
\definecolor{back-color}{rgb}{0.1686, 0.1686, 0.1686}
\definecolor{string-color}{rgb}{0.3333, 0.5254, 0.345}
\definecolor{key-color}{rgb}{0.9, 0.9, 0.9}
\definecolor{mygray}{rgb}{0.5,0.5,0.5}

\lstdefinestyle{my_vhdl_style}{
    language = vhdl,
    basicstyle = \ttfamily,
    otherkeywords = {;},
    keywordstyle = \color{blue},
    numbers=left,  % where to put the line-numbers; possible values are (none, left, right)
    %stepnumber=2,
    numbersep=5pt,                   % how far the line-numbers are from the code
    numberstyle=\tiny\color{mygray},
}

%.....Minitoc parameters.....
%\setcounter{minitocdepth}{2}
%\undotted

\title{BE : conception d'un accélérateur matériel\\pour le calcul de polynômes\\\textbf{listing}}
\author{Electronique numérique}

\begin{document}
\maketitle
\tableofcontents

\section{Controler}

\subsection{Machine d'états finis}
\lstinputlisting[style = my_vhdl_style]{./VHDL/controler.vhd}

\subsection{Package dédié}
\lstinputlisting[style = my_vhdl_style]{./VHDL/controler_pkg.vhd}

\section{Datapath}
\lstinputlisting[style=my_vhdl_style]{./VHDL/datapath.vhd}

\section{FSMD microarchitecture}
\lstinputlisting[style=my_vhdl_style]{./VHDL/fsmd.vhd}

\section{Testbench}
\lstinputlisting[style=my_vhdl_style]{./VHDL/fsmd_tb.vhd}

\section{Simulation}
\begin{center}
  \includegraphics[scale=0.3]{gtkwave_commente.png}
\end{center}

\end{document}
